\renewcommand{\Anno}{2004}

\renewcommand{\NumeroProblema}{3}

\renewcommand{\Testo}{Ad un torneo di tennis ad eliminazione diretta
partecipano 16 giocatori.  Il tabellone viene sorteggiato.
	
Si suppone che fra i giocatori esista una graduatoria di bravura nota
a priori, che viene sempre rispettata, nel senso che in ogni incontro
del torneo il giocatore migliore batte sempre quello peggiore (e non
esistono giocatori della stessa bravura).
	\begin{enumerate}
		\item  Determinare la probabilit\`{a} che tutti i 4 migliori 
		giocatori accedano alle semifinali.
	
		\item  Determinare la probabilit\`{a} che il giocatore sesto per 
		bravura acceda alle semifinali.
	\end{enumerate}
	
{\small Nota.  Il tabellone viene formato nel seguente modo: i nomi
dei 16 giocatori vengono tirati a sorte, senza tener conto della
bravura, e scritti di fila: a questo punto il primo sorteggiato sfida
il secondo sorteggiato, il terzo sfida il quarto e cos\`{i} via.  I
giocatori che hanno vinto questi primi 8 incontri si sfidano a loro
volta nei quarti di finale, nel senso che il giocatore che ha vinto il
primo incontro sfida il vincente del secondo incontro, e cos\`{i} via.
I 4 giocatori che vincono i quarti di finale si incontrano nelle
semifinali, sempre con lo stesso metodo (i vincitori dei primi due
quarti vanno alla prima semifinale).  Infine i vincitori delle
semifinali disputano la finale.}}

\renewcommand{\Risposta}{\begin{enumerate}
		\item  $64/455$ (circa il 14\%). 
	
		\item  $24/91$ (circa il 26\%).
	\end{enumerate}}

\renewcommand{\Aiutino}{Disegniamo il tabellone nel modo usuale, con i
sedici giocatori uno sotto l'altro in modo che i primi 4 producono il
primo semifinalista, quelli dal quinto all'ottavo il secondo
semifinalista, e cos\`{\i} via.

\begin{enumerate}
	\renewcommand{\labelenumi}{(\alph{enumi})} 
	\item In questo caso i
	4 migliori devono trovarsi in 4 quartine diverse.  Una volta
	piazzato il primo, quanti posti abbiamo a disposizione per i
	secondo?  Se va bene, quanti posti abbiamo per il terzo?  E per il
	quarto?

	\item Se il sesto non arriva in semifinale, vuol dire che nella
	sua quartina non c'\`{e} nessuno dei primi 5.  Scegliamo quindi il
	secondo giocatore da mettere nella quartina del sesto: in quanti
	modi possiamo farlo?  E il terzo?  E il quarto?
\end{enumerate}}

\renewcommand{\SoluzioneUfficiale}{\TitoloSoluzione Disegnando il 
tabellone nel modo usuale, i giocatori risultano suddivisi in 4 
quartine, ciascuna delle quali produce un semifinalista.
\begin{center}
	\psset{unit=1.5ex}
	\pspicture(0,0)(20,33)
	\multips(18.5,7)(0,16){2}{\psframe(0,0)(1,1)}
	\multips(24.5,15)(0,0){1}{\psframe(0,0)(1,1)}

	\psset{fillstyle=solid}

	\psset{fillcolor=cyan}
	\rput(0,0){\pspicture(0,0)
	\multips(0,0)(0,2){4}{\psframe(0,0)(1,1)}
	\psframe(6,1)(7,2)\psframe(6,5)(7,6)\psframe(12,3)(13,4)
	\endpspicture}
	
	\psset{fillcolor=green}
	\rput(0,8){\pspicture(0,0)
	\multips(0,0)(0,2){4}{\psframe(0,0)(1,1)}
	\psframe(6,1)(7,2)\psframe(6,5)(7,6)\psframe(12,3)(13,4)
	\endpspicture}
	
	\psset{fillcolor=red}
	\rput(0,16){\pspicture(0,0)
	\multips(0,0)(0,2){4}{\psframe(0,0)(1,1)}
	\psframe(6,1)(7,2)\psframe(6,5)(7,6)\psframe(12,3)(13,4)
	\endpspicture}
	
	\psset{fillcolor=yellow}
	\rput(0,24){\pspicture(0,0)
	\multips(0,0)(0,2){4}{\psframe(0,0)(1,1)}
	\psframe(6,1)(7,2)\psframe(6,5)(7,6)\psframe(12,3)(13,4)
	\endpspicture}
	
	\psset{fillstyle=none}

	\psset{arm=2.5\psunit}

	\multido{\n=0.5+4.0}{8}{\pcangle[angleB=180](1,\n)(!6 1 \n\space add)}
	\multido{\n=2.5+4.0}{8}{\pcangle[angleB=180](1,\n)(!6 1 neg \n\space add)}
	\multido{\n=1.5+8.0}{4}{\pcangle[angleB=180](7,\n)(!12 2 \n\space add)}
	\multido{\n=5.5+8.0}{4}{\pcangle[angleB=180](7,\n)(!12 2 neg \n\space add)}
	\pcangle[angleB=180](13,3.5)(18,7.5)
	\pcangle[angleB=180](13,19.5)(18,23.5)
	\pcangle[angleB=180](19,7.5)(24,15.5)
	\pcangle[angleB=180](13,11.5)(18,7.5)
	\pcangle[angleB=180](13,27.5)(18,23.5)
	\pcangle[angleB=180](19,23.5)(24,15.5)
	\endpspicture
\end{center}

\DomandaA Se i primi 4 per bravura accedono tutti alle semifinali, 
vuol dire che sono stati sorteggiati uno per quartina. Fissata la 
posizione del primo per bravura, abbiamo quindi 12 posizioni buone per il 
secondo in bravura sulle 15 disponibili. Fissato anche il secondo, 
rimangono 8 posizioni buone su 14 per il terzo in bravura ed infine 4 
posizioni buone su 13 per il quarto in bravura.

Ne segue che la probabilit\`{a} richiesta \`{e}
$$\frac{12}{15}\cdot\frac{8}{14}\cdot\frac{4}{13}=\frac{64}{455},$$
cio\`{e} circa del 14\%.
\DomandaB
Se il sesto per bravura arriva in semifinale, vuol dire che nella sua 
quartina non c'\`{e} nessuno dei primi 5. Possiamo dunque scegliere 
in 10 modi su 15 il secondo componente della sua quartina, poi in 9 modi 
su 14 il terzo, ed infine in 8 modi su 13 il quarto.

La probabilit\`{a} richiesta \`{e} pertanto
$$\frac{10}{15}\cdot\frac{9}{14}\cdot\frac{8}{13}=\frac{24}{91},$$
cio\`{e} circa del 26\%.}

\renewcommand{\SoluzioniAlternative}{\TitoloSoluzioneAlternativaA In 
entrambi i casi possiamo calcolare la probabilit\`{a} richiesta come 
rapporto tra casi favorevoli e casi possibili.

Per il punto (a), una possibilit\`{a} \`{e} di considerare le 
posizioni occupate nel tabellone dai 4 pi\`{u} bravi. I casi 
possibili, ovviamente equiprobabili, sono
$$\binom{16}{4}=1820,$$
e cio\`{e} tanti quanti i modi di scegliere 4 oggetti su 16. I casi 
favorevoli sono quelli in cui le 4 posizioni si trovano in quartine 
diverse, e sono quindi $4^{4}=256$, in quanto si tratta di scegliere 
una casella dalla prima quartina, una dalla seconda, una dalla terza 
ed una dalla quarta. Pertanto la probabilit\`{a} richiesta nel punto 
(a) \`{e}
$$\frac{256}{1820}=\frac{64}{455}.$$

Nel punto (b) possiamo considerare fissata la posizione del sesto in 
bravura e considerare le posizioni occupate dai 5 pi\`{u} bravi. I 
casi possibili, ovviamente equiprobabili, sono i modi di scegliere 5 
oggetti su 15, mentre quelli favorevoli sono i modi di scegliere 5 
oggetti su 12 (tanti sono i posti nelle 3 quartine che non contengono 
il sesto). Pertanto la probabilit\`{a} richiesta nel punto 
(b) \`{e}
$$\frac{\displaystyle{\binom{12}{5}}}{\displaystyle{\binom{15}{5}}}=
\frac{12\cdot 11\cdot 10\cdot 9\cdot 8}{5!}\cdot
\frac{5!}{15\cdot 14\cdot 13\cdot 12\cdot 11}=\frac{24}{91}.$$

In alternativa, sempre per il punto (b), possiamo considerare come
casi possibili i modi di scegliere i 3 ``compagni di quartina'' del
sesto in bravura (che sono i modi di scegliere 3 giocatori su 15) e
come casi favorevoli i modi di scegliere i 3 giocatori tra i 10 che
per bravura occupano i posti dal settimo in gi\`{u}.  In questo modo
la probabilit\`{a} si calcola come
$$\frac{\displaystyle{\binom{10}{3}}}{\displaystyle{\binom{15}{3}}}=
\frac{10\cdot 9\cdot 8}{3!}\cdot
\frac{3!}{15\cdot 14\cdot 13}=\frac{24}{91}.$$

Si noti che quest'ultimo \`{e} semplicemente un modo alternativo di
ottenere lo stesso calcolo fatto nella prima soluzione (tolti i $3!$,
il numeratore ed il denominatore sono gli stessi).}

\renewcommand{\Osservazioni}{
}

\renewcommand{\Commenti}{
}